\section{Department Pilot Delivery Status and Remaining Work}
\label{sec:department-pilot-remaining-delivery-plan}

\subsection{Purpose}

This section records the definitive delivery position of the AION Business Twin and
Department Pilot programme as at 20 July 2026.  It distinguishes the shared platform
and Finance capabilities that have already been delivered from the capability increments,
department designs, integrations and production controls that remain outstanding.

The target operating loop remains:

\begin{quote}
Function discovery $\rightarrow$ canonical function model $\rightarrow$ Boardroom context
$\rightarrow$ approved Boardroom package $\rightarrow$ Department Pilot task queue
$\rightarrow$ Pilot work and user conversation $\rightarrow$ artifacts and evidence
$\rightarrow$ Business Container writeback $\rightarrow$ Boardroom handback.
\end{quote}

\subsection{Completed Delivery Baseline}

The shared technical foundation and the first executable Finance vertical have been
completed.  The source is committed and synchronized with \texttt{origin/main} at commit
\texttt{b3f25716}, titled \texttt{add grounded Finance Pilot conversation runtime}.

The completed baseline comprises:

\begin{itemize}
  \item Business Foundation discovery and the canonical Business Container.
  \item The canonical Business Financial Model, Finance discovery workflow and durable
        Finance Department Intelligence.
  \item Read-only Xero connectivity and spreadsheet evidence upload, extraction, review
        and acceptance.
  \item The Products \& Services operating model, including offerings, delivery, labour,
        production, materials, stock, procurement, payment terms, overheads, debt and
        targets.
  \item The approved light Products \& Services interface, maintained in a modular source
        file without materially expanding \texttt{app.js}.
  \item The versioned \texttt{aion.department\_pilot.*.v1} contracts covering approved
        packages, tasks, progress, conversations, context references, approvals, artifacts,
        receipts, Department Intelligence patches and Boardroom handbacks.
  \item Backend-authoritative Department Pilot queues, conversations, audit history,
        guarded task transitions, stale-writer protection and restart recovery.
  \item Genuine live Boardroom approval and idempotent routing of Finance actions into the
        Finance Pilot queue and Central Pilot monitor.
  \item Permission-checked execution of approved read-only Finance analysis using the exact
        hash-bound canonical context approved by the founder.
  \item Finance report storage in the Business Container and registration under
        \texttt{Finance/Reports} in the File Cabinet.
  \item Completion receipts, Finance Intelligence writeback and durable result handback to
        the next Boardroom snapshot.
  \item A permanent Finance Pilot conversation session grounded in selectively retrieved
        canonical evidence, with reliability, source, revision and content-hash labels.
  \item Period-aware management-account reporting covering profit and loss, cash and working
        capital; supported calculated margins and cost relationships; actual-versus-budget
        and previous-period comparisons; evidence freshness; verification state; and visible
        missing or conflicting evidence.
  \item Idempotent routing and execution, visible completion and failure states, deterministic
        grounded fallback behaviour and explicit design gates for unfinished departments.
  \item A verification baseline of 56 focused current-architecture automated tests, JavaScript and Python syntax
        checks, and successful packaged-desktop verification.
\end{itemize}

The Central Pilot remains the cross-department coordinator and monitor.  It does not replace
specialist Department Pilots.  Finance is currently the only specialist enabled to receive
canonical Boardroom work.  Sales, Marketing, Support, HR and Operations remain intentionally
design-gated.

\subsection{Remaining Delivery Order}

The following work is outstanding and shall be delivered in the stated order.

\subsubsection{Finance Management Reporting---Completed 20 July 2026}

The richer management-account reporting layer has been delivered.  It now:

\begin{itemize}
  \item reports revenue, direct costs, gross profit, gross-margin percentage, overheads,
        operating profit, supported net profit, cash movement, debtors and creditors;
  \item compares actual results with canonical budgets and compares the current period with the previous
        equivalent period;
  \item supports clearly identified reporting periods when that period is evidenced;
  \item displays source dates, evidence freshness, verification state and all missing, stale,
        assumed or unverified values;
  \item produces a readable Boardroom-ready management-account presentation in the Pilot handback;
  \item consumes accepted spreadsheet facts while withholding raw external payloads;
  \item refuses to silently choose between conflicting accepted facts; and
  \item stores generated reports and their receipts under \texttt{Finance/Reports}.
\end{itemize}

\subsubsection{Finance Forecasting and Scenario Modelling}

Finance shall subsequently support short-term cash-flow forecasting and controlled scenario
analysis for revenue, sales volume, price, discounts, costs, hiring, labour, stock purchases,
procurement and cash.  Best-case, expected-case and downside scenarios shall identify their
assumptions separately from verified historical facts.  Scenario results shall be available
to both the Finance Pilot and the Boardroom.

\subsubsection{Xero and Spreadsheet Reconciliation}

A reconciliation layer shall compare accepted spreadsheet facts with Xero for equivalent
reporting periods.  It shall detect conflicts, duplicates, missing accounts, stale sources
and inconsistent dates; identify the current authoritative source; and require visible user
review before accepted canonical facts are replaced.  Reconciliation decisions shall retain
full provenance.  Xero shall remain read-only during this increment.

\subsubsection{Finance Proposed Actions and Exact Approval}

Finance shall be able to stage proposed actions without executing them silently.  Every
proposal shall include the exact payload, expected financial effect, supporting evidence,
required permission and approval state.  Approval, rejection, amendment, execution and
failure shall generate immutable receipts.  No external financial write shall occur without
an exact, inspectable approval enforced by the backend.

\subsubsection{Recurring Finance Work}

Backend-authoritative schedules shall be added for weekly cash reporting, overdue-invoice
monitoring, monthly management accounts, period-close checks, budget-variance alerts,
cash-buffer warnings and recurring Boardroom updates.  Each schedule shall retain its last
run, next run and failure state.  Read-only work may run automatically; any external side
effect shall require a separate approval boundary.

\subsubsection{Finance Production Hardening}

Before Finance is declared production-complete, the system shall add task retry,
cancellation and recovery controls; improve failure diagnostics and provider timeout
handling; enforce bounded retrieval for large datasets; and complete business-isolation,
permission, sensitive-data and packaged-application end-to-end testing.  Reports, receipts,
Business Container records and File Cabinet pointers must remain valid after restart.

The Finance exit gate is the proven loop:

\begin{quote}
Boardroom decision $\rightarrow$ Finance Pilot queue $\rightarrow$ grounded analysis and
conversation $\rightarrow$ report and evidence $\rightarrow$ Finance Intelligence
$\rightarrow$ next Boardroom meeting.
\end{quote}

\subsubsection{Sales Design and Implementation}

Sales shall be designed with the founder before \texttt{sales\_pilot} is activated.  The
design shall cover Sales discovery, customers, prospects, pipeline stages, opportunity value,
CRM and ecommerce sources, inbound qualification, outbound calling, human handoff, email
funnels, consent, quotes, proposals, follow-up, forecasting and payment terms.  Revenue and
customer-payment implications shall reference Products \& Services and Finance canonically.

\subsubsection{Marketing Design and Implementation}

The Marketing design shall define audiences, segments, offers, channels, campaigns, budgets,
lead generation, attribution, the Marketing-to-Sales handoff, content preparation and
performance reporting.  Existing Marketing work shall be consolidated into the shared
Department Pilot runtime.  Publishing and advertising expenditure shall remain approval-gated.

\subsubsection{Support Design and Implementation}

Support shall define telephone, email and messaging capabilities; decide which communication
channels are native and which are plugins; and model customer history, issue categories,
service levels, knowledge, drafted responses, escalation and human handoff.  Outbound
communication shall be controlled.  Product issues, retention signals and cost implications
shall be shared with the relevant canonical departments.

\subsubsection{HR and People Design and Implementation}

HR shall include a visual organisation chart and accountability canvas covering founders,
employees, contractors, vacancies, reporting lines, responsibilities, skills, capacity and
employment cost.  Application roles and data/action permissions shall remain separate from
job titles.  Sensitive people information shall use stricter retrieval and visibility rules.
The resulting people and capacity context shall be available safely to Finance and Operations.

\subsubsection{Operations Design and Implementation}

Operations shall be designed last because it consumes the canonical models established by
the other functions.  Its scope shall cover delivery workflows, work queues, capacity,
production, labour, materials, stock, work in progress, procurement, suppliers, lead times,
bottlenecks, quality and operational performance.  Operations shall use shared canonical
references and shall not become a second Central Pilot.

\subsubsection{Connector Expansion}

Connector delivery shall be prioritised by demonstrated user demand.  Expected candidates
include QuickBooks and Sage; CRM platforms; Shopify, WooCommerce, Magento and Square;
payment processors; inventory and warehouse systems; supplier and purchasing systems; and
communications and support platforms.  Each connector requires authentication, canonical
mapping, sync state, provenance, visible errors and explicit approval boundaries.

\subsubsection{Catalogue and Transaction Scale}

Large catalogues and high-volume transaction histories shall move from monolithic JSON
documents to a canonical indexed store.  The Business Container shall retain compact summary
projections while Pilots retrieve only the records required for the current question or task.
The implementation shall support pagination, search, filtering and import reconciliation and
shall prevent entire catalogues or ledgers from being inserted into every AI prompt.

\subsection{Definition of Full Production Completion}

The full Department Pilot programme is production-complete only when:

\begin{itemize}
  \item every material value retains provenance, reporting period and verification state;
  \item department tasks and conversations survive restart and authenticated sign-in;
  \item approved Boardroom work routes reliably to the correct specialist Pilot;
  \item permissions are enforced by the backend rather than only by the interface;
  \item every external side effect requires an exact, inspectable approval;
  \item every attempted execution creates a success, failure or blocked receipt;
  \item artifacts are accessible from the correct Business Container and File Cabinet path;
  \item Department Intelligence and Boardroom context update without duplicating raw evidence;
  \item cross-department canonical references remain consistent; and
  \item automated and packaged-desktop end-to-end tests cover the complete operating loop.
\end{itemize}

\subsubsection{Completed Finance Forecasting and Scenario Slice}

The Finance Pilot now provides a deterministic, read-only planning model covering a horizon
of one to 36 months. It derives its monthly baseline from the canonical management-account
record and refuses to run when revenue, direct costs, overheads or opening cash are missing.
The founder may model annual revenue growth; revenue, price and volume changes; direct-cost
and overhead changes; monthly hiring cost; a one-off stock purchase; and a protected minimum
cash reserve. Each run produces expected, upside and downside monthly projections for
revenue, gross profit, operating profit, net cash movement and closing cash, including the
months in which the protected reserve would be breached.

Scenario inputs and results are stored as a separate hash-bound planning record under
\texttt{Finance/Scenarios}. The File Cabinet contains an index pointer, Finance Department
Intelligence receives a compact scenario summary, and the Boardroom snapshot receives the
latest variant summaries. Every record is labelled \texttt{modelled\_not\_actual}; running a
scenario does not change the Business Financial Model, accepted spreadsheet facts, Xero
evidence or historical accounts. The calculation policy and excluded movements are recorded
with the artifact so the Boardroom cannot mistake the model for verified accounting truth.

\subsubsection{Completed Xero-to-Spreadsheet Reconciliation Slice}

The Finance Pilot now compares accepted spreadsheet facts with the latest read-only Xero
report evidence. A comparison is made only where the reporting periods match. Each metric is
classified as matching, conflicting, missing from the spreadsheet, missing from Xero,
missing a reporting period or belonging to a different period. The resulting reconciliation
is hash-bound, stored under \texttt{Finance/Reconciliations}, indexed in the File Cabinet and
summarised into Finance Department Intelligence and the Boardroom snapshot.

The founder may record a decision to keep the spreadsheet, prefer Xero, investigate further
or make no change. This review does not itself replace canonical financial values. A decision
that would change accepted truth remains pending a separate exact-payload canonical update,
thereby preventing a reconciliation click from silently rewriting historical accounts. Xero
remains read-only throughout.

\subsubsection{Completed Exact-Payload Finance Proposal Slice}

The Finance Pilot now stages only explicitly permitted draft actions through the canonical
proposed-tool-call contract.  Each proposal displays its complete JSON payload and immutable
payload hash.  A positive decision is rejected unless the deciding user supplies that exact
hash; consequently an approval cannot authorise an amended amount, recipient, account,
provider, task, department or business.  The decision is persisted inside the canonical work
envelope and is visible in both the Finance and Central Pilot queues.

Approval does not itself enable execution.  The runtime performs a second permission check
against the Finance Pilot's live-connector allow-list.  That allow-list remains empty in this
increment, so an approved proposal reports
\texttt{live\_connector\_permission\_not\_enabled} and records that no external action was
performed.  This proves the approval boundary without granting Xero or any other provider a
write capability.

\subsubsection{Completed Recurring Finance and Recovery Slice}

The backend now persists four standard read-only schedules: weekly cash updates, monthly
management reports, monthly period-close checks and daily exception or cash-buffer alerts.
Each hash-bound schedule records its cadence, enabled state, next and last run, last outcome,
run count, failure count and latest diagnostic.  Every successful run creates a separate
hash-bound artifact in the Finance Business Container, an index pointer under
\texttt{Finance/Recurring}, a compact Finance Intelligence projection and a latest-run
reference for the Boardroom.  Scheduled results explicitly record that no external action
was performed.

Operators may pause, resume and run schedules manually.  Failed Department Pilot tasks now
expose durable diagnostics and may be retried within their configured retry budget, recovered
to the specialist queue or cancelled.  Recovery does not replay a prior connector action.
Business and department identity are checked before diagnostics or state changes are returned.

\subsubsection{Completed Finance Finalisation and Installed-Application Validation}

The Finance review screen now finalises the first Business Financial Model through a
canonical backend persistence transaction.  Uploaded workbook analysis is represented in
the local model by compact, hash-bound references rather than being copied wholesale into
browser storage.  This removes the storage-quota failure that previously caused the
\texttt{Build Financial Model} control to appear unresponsive.  The interface displays a
building state, requires a canonical Business Container receipt before completion, and
returns a visible recoverable error if finalisation fails.

The packaged application has been exercised across a complete restart.  The validation
proved persistence and retrieval of an approved Finance Boardroom action, its specialist
queue state, an exact-payload approval decision, a recurring cash schedule, its completed
run state and the resulting File Cabinet pointer.  The approved action remained blocked
from external execution because no live write permission was enabled.  Scenario,
reconciliation and recurring-work routes were also tested after removing generic-route
shadowing.  The installed executable passed when started directly; normal macOS Finder or
\texttt{open} launch did not remain alive reliably in the test environment.  Launch-process
ownership therefore remains an explicit packaging exit-gate item.

\subsubsection{Finance Permission and Sensitive-Data Matrix}

The backend now publishes a canonical Finance security policy.  The governing rules are:

\begin{itemize}
  \item external financial writes are denied by default;
  \item approval is bound to the immutable hash of the exact proposed payload;
  \item connector credentials and refresh tokens remain in macOS Keychain and are never
        written to the Business Container, Pilot conversation or Boardroom prompt;
  \item financial summaries may be used for Finance and Boardroom analysis;
  \item transaction-level details are retrieved only when required for a bounded task;
  \item counterparty personal data is masked or minimised outside the authorised Finance task;
  \item payroll and tax-sensitive data is restricted from general Boardroom projection;
  \item draft invoices, reminders or credit-control actions may be prepared but not sent;
  \item accounting reconciliation posting, journal posting, payments, transfers and tax filing
        remain disabled; and
  \item every connector attempt must produce a secret-free success, retry, failure or blocked receipt.
\end{itemize}

Xero connector attempts now persist such receipts under
\texttt{Finance/Connector Receipts}.  A receipt records provider, operation, attempt, status,
provider status code, retryability, error classification, request hash and the invariant
\texttt{external\_write\_performed=false}; it excludes access tokens, request headers and raw
provider error bodies.

\subsubsection{Current Finance Capability Boundary}

Finance is presently a read-only financial analyst, modeller and controlled proposal system.
It can produce management reports and KPIs; model expected, upside and downside scenarios;
monitor cash and recurring exceptions; compare aligned spreadsheet and Xero report summaries;
answer grounded questions; and prepare exact-payload proposals.  It does not yet perform
transaction-level bookkeeping.  In particular, it cannot match individual bank entries to
invoices or bills, post a reconciliation or journal, amend Xero, send collection messages,
issue an invoice, transfer money or file a tax return.

\subsubsection{Next Finance Director Transaction and Decision Vertical}

Before Sales implementation begins, Finance shall be completed as a useful operating Finance
Director through the following bounded increments:

\begin{enumerate}
  \item Create a canonical ledger and transaction layer for chart-of-account entries, bank
        transactions, invoices, bills, payments, journals, projects and tracking categories.
  \item Implement transaction-level matching with confidence scores, duplicate and missing
        record detection, exception queues and founder or accountant review.
  \item Add an inspectable reconciliation workspace.  Suggested matches shall remain drafts;
        posting requires an individually enabled connector capability and exact-payload approval.
  \item Extend the living Business Financial Model across profit and loss, balance sheet, cash
        flow, receivables, payables, working capital, budgets, job or project profitability,
        work in progress, pipeline and Products \& Services unit economics.
  \item Add proactive signals for cash runway, protected reserves, margin erosion, overdue
        debtors, quote leakage, work in progress, capacity and pipeline risk.
  \item Translate analysis into canonical Boardroom proposals and cross-department requests.
        Finance may recommend pricing changes, invoice chasing, quote follow-up or demand
        generation, while Sales, Marketing and Operations retain their execution permissions.
  \item Introduce narrow write adapters only after draft, review, approval, execution receipt
        and rollback or recovery behaviour has been proved.  Generic Finance write permission
        shall never be granted.
\end{enumerate}

The Finance exit gate is reached when the Pilot can explain the current financial position,
forecast the consequence of a proposed decision, identify the transactional or operational
drivers of a risk, prepare an inspectable response, obtain the precise approval required and
return verified results to the Business Container and the next Boardroom meeting.  Sales
capability design begins after this gate.

\subsubsection{Completed Finance Director Transaction and Decision Vertical}

The read-only Xero integration now feeds a canonical transaction ledger covering accounts,
sales invoices, supplier bills, bank transactions, payments, manual journals, tracking
categories and projects.  The source snapshot is hash checked before use.  Large collections
are stored as bounded JSON Lines records below \texttt{Finance/Ledgers}; the Financial Model,
Department Intelligence, Boardroom and File Cabinet receive only compact summaries and
canonical references.

The transaction reconciliation service proposes payment-to-invoice and bank-entry-to-invoice
matches with explicit confidence scores and reasons.  It detects duplicate records, unmatched
entries, overdue receivables and overdue payables and presents them in an exception queue.  An
authorised reviewer may accept a match as draft intent, reject it or require investigation.
None of those decisions posts a reconciliation, journal or other mutation to Xero.

The Finance Director model combines the canonical ledger, accepted Business Financial Model
and Products \& Services operating model.  It produces an evidence-labelled operating profit
and loss view, partial balance sheet, cash-flow proxy, working-capital position, project,
budget, work-in-progress, pipeline and unit-economics context.  Cash runway is based on net
monthly cash burn rather than gross monthly expenditure.  Proactive signals cover protected
reserve breaches, runway, low margin, operating losses, overdue invoices, quote leakage, weak
pipeline, cash tied up in work in progress, absent unit economics and unverified evidence.

Each supported signal may create an inspectable Boardroom-review proposal for invoice chasing,
quote follow-up, pricing scenarios, controllable-cost review, Sales pipeline work, Marketing
demand-generation options or Operations work-in-progress completion.  Proposal creation does
not route, send or execute an action.  The permanent Finance terminal exposes the living model,
ledger controls, warnings, proposals, match review and exception queue.

The backend permission matrix permits read-only ledger construction, transaction matching and
reviewed draft intent.  Reconciliation posting, journal posting, payments, transfers and tax
filing remain prohibited.  A future provider-specific write adapter must be separately enabled,
bound to an exact approved payload, produce immutable success or failure receipts and prove
recovery behaviour.  Generic Finance or Xero write permission shall never be granted.

\subsubsection{Reconciliation Intelligence and Xero Provider Boundary}

Counterparty classification uses a hybrid decision model.  A mapping already confirmed by an
authorised reviewer is reused deterministically.  An account code already recorded on the Xero
transaction remains provider evidence.  Only an unknown counterparty is eligible for bounded
LLM assistance.  The model receives a minimised description and the existing chart of accounts;
it may suggest an existing code but cannot invent a canonical account, confirm its own mapping
or perform an external write.  Invalid account codes and malformed confidence values are treated
as unsupported evidence.  New mappings remain pending until a person confirms them, and uncertain
cases become explicit clarification requests.  LLM classification is capped per reconciliation
run to limit disclosure, latency and cost.

Xero's Accounting API does not expose bank-statement-line reconciliation.  AION therefore does
not simulate such a write or describe an accepted local match as a completed Xero reconciliation.
For an accepted draft match it can prepare a provider-specific manual handoff containing the
exact transaction, invoice or payment references.  Approval is bound to the immutable payload
hash.  Execution produces an idempotent, hash-verified provider-limitation receipt stating that
the action must be completed in Xero Bank Rec or Cash Coding, that zero provider reconciliations
were posted and that \texttt{external\_write\_performed=false}.  The receipt is indexed under
\texttt{Finance/Xero Reconciliation Handoffs} and is isolated to the canonical Business Container.

The Finance interface also compares granted Xero scopes with the current read-only requirement.
An existing connection missing manual-journal or project scopes presents a clear
\emph{Reauthorise Xero} action before synchronisation.  Supported future Xero mutations, such as
payment creation, shall be separate provider capabilities with their own payload, permission,
approval and receipt contracts; they shall never be used to imply bank-statement reconciliation.
